\documentclass[11pt]{article}
\usepackage{amsmath,amssymb,amsthm,mathtools}
\usepackage[margin=1in]{geometry}
\usepackage{hyperref}

\title{One-Step Full-Batch Gradient Descent Scaling in a Tied-Embedding Toy Transformer with Zipf Data}
\author{}
\date{}

\begin{document}
\maketitle

\section{Model and data}

Let the vocabulary size be $V$ and hidden width be $n$.
Let $E\in\mathbb{R}^{n\times V}$ be an embedding matrix (tied for input and output)
and $W\in\mathbb{R}^{n\times n}$ be a hidden linear map.
For a token $t\in\{1,\dots,V\}$ with one-hot vector $e_t\in\mathbb{R}^{V}$, define
\[
x^{(t)} := E e_t\in\mathbb{R}^n,\qquad
h^{(t)} := W x^{(t)}\in\mathbb{R}^n,\qquad
z^{(t)} := \frac{1}{n}E^\top h^{(t)}\in\mathbb{R}^V.
\]
The prediction is $p^{(t)} := \mathrm{softmax}(z^{(t)})$ and the per-example loss is the cross-entropy
\[
\ell^{(t)} := -\log p^{(t)}_t.
\]

\paragraph{Initialization.}
Assume iid Gaussian initialization
\[
W_{ij}\sim \mathcal{N}\!\left(0,\frac{1}{n}\right)\quad \text{iid},\qquad
E_{ik}\sim \mathcal{N}(0,1)\quad \text{iid},
\]
with $W$ independent of $E$.

\paragraph{Data distribution (Zipf).}
Assume full-batch (population) gradient descent under a token distribution
$q=(q_1,\dots,q_V)$, where $t\sim q$ and the label is $y=t$.
For Zipf with exponent $s>0$ (tokens ranked in decreasing frequency),
\[
q_k = \frac{k^{-s}}{H_{V,s}},\qquad H_{V,s}:=\sum_{r=1}^V r^{-s}.
\]
The population loss is
\[
\mathcal{L} := \mathbb{E}_{t\sim q}\big[\ell^{(t)}\big] = \sum_{t=1}^V q_t\,\ell^{(t)}.
\]

\section{Forward-pass scales at initialization}

\subsection{Embedding coordinates}
For fixed $t$, the embedding vector is the $t$-th column of $E$:
\[
x^{(t)}_i = E_{it}\sim \mathcal{N}(0,1)\quad \text{exactly},
\]
so $\|x^{(t)}\|^2=\sum_{i=1}^n (x^{(t)}_i)^2\sim \chi^2_n$ and therefore
\[
\frac{1}{n}\|x^{(t)}\|^2 \to 1 \quad \text{in probability as } n\to\infty.
\]

\subsection{Hidden coordinates}
Condition on $x^{(t)}$. Each row of $W$ is independent of $x^{(t)}$, hence for each $i$,
\[
h^{(t)}_i = \sum_{j=1}^n W_{ij}x^{(t)}_j \,\Big|\, x^{(t)}
\sim \mathcal{N}\!\left(0,\frac{1}{n}\|x^{(t)}\|^2\right),
\]
and $\{h^{(t)}_i\}_{i=1}^n$ are conditionally iid. In particular, since $\|x^{(t)}\|^2/n\to 1$,
each coordinate $h^{(t)}_i$ is $\Theta(1)$ and $\|h^{(t)}\|^2=\Theta(n)$ with high probability (w.h.p.).

\subsection{Logits}
For $k\neq t$, conditional on $h^{(t)}$ the column $E_{\cdot k}$ is independent of $h^{(t)}$, so
\[
z^{(t)}_k = \frac{1}{n}\sum_{i=1}^n E_{ik}h^{(t)}_i \,\Big|\, h^{(t)}
\sim \mathcal{N}\!\left(0,\frac{1}{n^2}\|h^{(t)}\|^2\right).
\]
Since $\|h^{(t)}\|^2=\Theta(n)$, it follows that
\[
z^{(t)}_k = \Theta\!\left(\frac{1}{\sqrt{n}}\right)\quad \text{w.h.p. for each fixed }k\neq t.
\]
For $k=t$ (tied column), $z^{(t)}_t = \frac{1}{n}(x^{(t)})^\top W x^{(t)}$ is also
$\Theta(n^{-1/2})$ w.h.p. (as a Gaussian quadratic form conditional on $x^{(t)}$).

Thus all logits are near $0$ at init, with typical coordinate size $\Theta(n^{-1/2})$, regardless of $V$.

\section{Cross-entropy logit gradients and the effect of large $V$}

For an example $t$, the logit-gradient is
\[
g^{(t)} := \nabla_{z^{(t)}}\ell^{(t)} = p^{(t)} - e_t.
\]
Since $z^{(t)}=\Theta(n^{-1/2})$ at init, $p^{(t)}$ is near-uniform:
\[
p^{(t)}_k = \frac{1}{V} + o(1)\quad \text{(for each fixed }k\text{)},
\]
hence to leading order
\[
g^{(t)} \approx u - e_t,\qquad u:=\frac{1}{V}\mathbf{1}.
\]
In particular,
\[
g^{(t)}_t \approx -1,\qquad g^{(t)}_{k\neq t}\approx \frac{1}{V}.
\]
Although most coordinates are $\Theta(1/V)$, the $\ell_2$ norm remains constant-order:
\[
\|g^{(t)}\|_2^2 = (1-p^{(t)}_t)^2 + \sum_{k\neq t} (p^{(t)}_k)^2 = \Theta(1).
\]

\paragraph{Full-batch (population) average.}
Define the average logit-gradient under $t\sim q$:
\[
\bar g := \mathbb{E}_{t\sim q}[g^{(t)}] = \sum_{t=1}^V q_t g^{(t)}.
\]
Using $g^{(t)}\approx u-e_t$ gives
\[
\bar g \approx \sum_{t=1}^V q_t(u-e_t) = u-q,
\]
so coordinatewise
\[
\bar g_k \approx \frac{1}{V} - q_k.
\]
Under Zipf, $q_k = k^{-s}/H_{V,s}$, this exhibits a head/tail split:
if $q_k\gg 1/V$, then $\bar g_k\approx -q_k$; if $q_k\ll 1/V$, then $\bar g_k\approx 1/V$.

The rank cutoff $k_\star$ such that $q_{k_\star}\asymp 1/V$ satisfies
\[
\frac{k_\star^{-s}}{H_{V,s}} \asymp \frac{1}{V}
\quad \Longrightarrow \quad
k_\star \asymp \left(\frac{V}{H_{V,s}}\right)^{1/s}.
\]
For $s>1$, $H_{V,s}\to \zeta(s)$ so $k_\star=\Theta(V^{1/s})$; for $s=1$, $H_{V,1}\sim \log V$ so
$k_\star=\Theta(V/\log V)$.

\section{Backpropagation for one example}

For a fixed example $t$, define
\[
\delta_h^{(t)} := \nabla_{h^{(t)}}\ell^{(t)},\qquad
\delta_x^{(t)} := \nabla_{x^{(t)}}\ell^{(t)}.
\]
Using $z^{(t)}=\frac{1}{n}E^\top h^{(t)}$ and $h^{(t)}=Wx^{(t)}$,
\[
\delta_h^{(t)} = \frac{1}{n}E g^{(t)},\qquad
\delta_x^{(t)} = W^\top \delta_h^{(t)}.
\]
At init, since $\|g^{(t)}\|_2=\Theta(1)$, each coordinate of $\delta_h^{(t)}$ has typical size $\Theta(1/n)$,
and each coordinate of $\delta_x^{(t)}$ also has typical size $\Theta(1/n)$.

\section{Full-batch gradients and one GD step}

Population gradients:
\[
\nabla_W \mathcal{L} = \sum_{r=1}^V q_r\,\delta_h^{(r)} (x^{(r)})^\top,
\]
\[
\nabla_E \mathcal{L} = \sum_{r=1}^V q_r\left(\delta_x^{(r)} e_r^\top + \frac{1}{n}h^{(r)} (g^{(r)})^\top\right).
\]
Take one full-batch GD step with group learning rates $\eta_W,\eta_E$:
\[
\Delta W := W^+ - W = -\eta_W \nabla_W\mathcal{L},\qquad
\Delta E := E^+ - E = -\eta_E \nabla_E\mathcal{L}.
\]

\subsection{Induced updates in activations (first order)}
For a fixed token $t$, define updated activations using $(W^+,E^+)$:
\[
x^{(t)+} := E^+ e_t,\qquad h^{(t)+}:=W^+ x^{(t)+},\qquad z^{(t)+} := \frac{1}{n}(E^+)^\top h^{(t)+}.
\]
To first order in $(\eta_W,\eta_E)$ (dropping terms quadratic in step sizes),
\[
\Delta x^{(t)} := x^{(t)+}-x^{(t)} = (\Delta E)e_t,
\]
\[
\Delta h^{(t)} := h^{(t)+}-h^{(t)} \approx (\Delta W)x^{(t)} + W(\Delta x^{(t)}),
\]
\[
\Delta z^{(t)} := z^{(t)+}-z^{(t)} \approx \frac{1}{n}\Big((\Delta E)^\top h^{(t)} + E^\top \Delta h^{(t)}\Big).
\]

\section{Where Zipf enters: per-token magnitude of full-batch updates}

Because the population gradients are weighted by $q_r$, the leading contribution to the update experienced
by token $t$ is proportional to $q_t$ (the frequency of $t$ under $q$). At init, the relevant
per-coordinate backprop signals satisfy $\delta_h^{(t)}=\Theta(1/n)$ and $\delta_x^{(t)}=\Theta(1/n)$,
hence the typical per-coordinate update magnitudes obey the following scaling laws:

\subsection{Embedding update magnitude}
The input-side part of $\nabla_E\mathcal{L}$ contributes to column $t$ in proportion to $q_t$,
yielding
\[
\Delta x^{(t)} = (\Delta E)e_t \quad \Longrightarrow \quad
|\Delta x^{(t)}_i| = \Theta\!\left(\eta_E \frac{q_t}{n}\right).
\]
Under Zipf, $q_t = t^{-s}/H_{V,s}$, so
\[
|\Delta x^{(t)}_i| = \Theta\!\left(\eta_E \frac{1}{n}\frac{t^{-s}}{H_{V,s}}\right).
\]

\subsection{Hidden update magnitude}
Using $\Delta h^{(t)}\approx (\Delta W)x^{(t)} + W\Delta x^{(t)}$ and the rank-1 structure
$\nabla_W \ell^{(t)}=\delta_h^{(t)}(x^{(t)})^\top$, the dominant \emph{self} contribution to $(\Delta W)x^{(t)}$
scales like $\eta_W q_t$ per coordinate at init, while $W\Delta x^{(t)}$ contributes at most on the order of
$\eta_E q_t/n$ per coordinate. Thus
\[
|\Delta h^{(t)}_i| = \Theta(\eta_W q_t) + \Theta\!\left(\eta_E \frac{q_t}{n}\right),
\]
and typically the $\Theta(\eta_W q_t)$ term dominates unless $\eta_E\gg n\eta_W$.

Under Zipf,
\[
|\Delta h^{(t)}_i| = \Theta\!\left(\eta_W \frac{t^{-s}}{H_{V,s}}\right)
\quad \text{(dominant term)}.
\]

\subsection{Logit update magnitude and the $V$ vs $\sqrt{n}$ regime split}

A key object is the Gram matrix $G:=E^\top E\in\mathbb{R}^{V\times V}$.
At init, for $j\neq k$ one has $G_{jk}=\sum_{i=1}^n E_{ij}E_{ik}=\Theta_p(\sqrt{n})$,
so $(1/n)G_{jk}=\Theta_p(1/\sqrt{n})$, while $(1/n)G_{kk}\approx 1$.

The dominant self contribution to $\Delta z^{(t)}$ is proportional to
\[
q_t\left(\frac{1}{n}E^\top E\right)g^{(t)}.
\]
Since $g^{(t)}\approx u-e_t$, the label coordinate satisfies
\[
\left[\left(\frac{1}{n}E^\top E\right)g^{(t)}\right]_t = -1 + O_p\!\left(\frac{1}{\sqrt{n}}\right),
\]
whereas for a non-label coordinate $j\neq t$ one gets
\[
\left[\left(\frac{1}{n}E^\top E\right)g^{(t)}\right]_j
= \frac{1}{V} + \Theta_p\!\left(\frac{1}{\sqrt{n}}\right).
\]
Therefore the non-label behavior has two regimes:
\[
\frac{1}{V} \gg \frac{1}{\sqrt{n}} \ \ (V\ll \sqrt{n})\quad \Longrightarrow\quad
\text{non-label scale } \Theta(1/V),
\]
\[
\frac{1}{\sqrt{n}} \gg \frac{1}{V} \ \ (V\gg \sqrt{n})\quad \Longrightarrow\quad
\text{non-label scale } \Theta(1/\sqrt{n}).
\]

Combining with the full-batch frequency factor $q_t$ and the effective prefactor
$\eta_W+\eta_E/n$ that appears in the induced logit change, one obtains
\[
|\Delta z^{(t)}_t|
= \Theta\!\Big(\big(\eta_W+\eta_E/n\big)\,q_t\Big),
\]
\[
|\Delta z^{(t)}_j| \ (j\neq t)
= \Theta\!\Big(\big(\eta_W+\eta_E/n\big)\,q_t\cdot \max\{1/V,\ 1/\sqrt{n}\}\Big).
\]
Under Zipf $q_t=t^{-s}/H_{V,s}$, these become
\[
|\Delta z^{(t)}_t|
= \Theta\!\Big(\big(\eta_W+\eta_E/n\big)\frac{t^{-s}}{H_{V,s}}\Big),
\]
\[
|\Delta z^{(t)}_j| \ (j\neq t)
= \Theta\!\Big(\big(\eta_W+\eta_E/n\big)\frac{t^{-s}}{H_{V,s}}\cdot \max\{1/V,\ 1/\sqrt{n}\}\Big).
\]

\section{Learning-rate normalization choices under Zipf}

Because $q_t$ varies widely, no single global learning rate can make one-step updates
$\Theta(1)$ simultaneously for all token ranks $t$. Two natural normalizations are:

\subsection{Head-token normalization ($t=1$)}
Since $q_1 = 1/H_{V,s}$, achieving $\Theta(1)$ updates for the most frequent token suggests
\[
\eta_W = \Theta(H_{V,s}),\qquad \eta_E = \Theta(n H_{V,s}),
\]
which for $s>1$ implies $\eta_W=\Theta(1)$, $\eta_E=\Theta(n)$, while for $s=1$ implies
$\eta_W=\Theta(\log V)$, $\eta_E=\Theta(n\log V)$.

\subsection{Typical-draw normalization ($t\sim q$)}
A frequency scale for a random draw is $\mathbb{E}_{t\sim q}[q_t]=\sum_{k=1}^V q_k^2=\|q\|_2^2$.
Under Zipf,
\[
\|q\|_2^2 = \frac{H_{V,2s}}{H_{V,s}^2}.
\]
To keep typical per-token updates $\Theta(1)$ in expectation over $t\sim q$, one would scale
\[
\eta_W = \Theta\!\left(\frac{1}{\|q\|_2^2}\right)=\Theta\!\left(\frac{H_{V,s}^2}{H_{V,2s}}\right),
\qquad
\eta_E = \Theta\!\left(\frac{n}{\|q\|_2^2}\right)=\Theta\!\left(n\frac{H_{V,s}^2}{H_{V,2s}}\right).
\]
For $s=1$, $H_{V,1}\sim \log V$ and $H_{V,2}\to \zeta(2)$, hence
\[
\eta_W=\Theta((\log V)^2),\qquad \eta_E=\Theta(n(\log V)^2).
\]
For $s>1$, $H_{V,s},H_{V,2s}$ converge, yielding again $\eta_W=\Theta(1)$, $\eta_E=\Theta(n)$.

\section{Summary}

At initialization, $x^{(t)}_i=\Theta(1)$, $h^{(t)}_i=\Theta(1)$, and $z^{(t)}_k=\Theta(n^{-1/2})$ w.h.p.
independent of $V$.

Under full-batch GD with Zipf token frequencies, one-step update magnitudes for a rank-$t$ token
scale proportionally to $q_t=t^{-s}/H_{V,s}$:
\[
|\Delta x^{(t)}_i|=\Theta\!\left(\eta_E \frac{q_t}{n}\right),\qquad
|\Delta h^{(t)}_i|=\Theta(\eta_W q_t)\ \text{(dominant)},
\]
\[
|\Delta z^{(t)}_t|=\Theta\!\Big((\eta_W+\eta_E/n)q_t\Big),
\]
and for non-label logits $j\neq t$ there is a regime split in the tied model:
\[
|\Delta z^{(t)}_j|=\Theta\!\Big((\eta_W+\eta_E/n)q_t\cdot \max\{1/V,\ 1/\sqrt{n}\}\Big).
\]
Thus when $V\gg \sqrt{n}$, non-label coordinate changes are typically governed by the
Gram off-diagonal scale $\Theta(1/\sqrt{n})$ rather than $1/V$.

\section{Goal and scope}

We study the \emph{one-step} optimal learning rate $\eta^\star(\beta,n)$ for two update rules:
(i) SGD and (ii) Muon (polar update), in a softmax model whose representation covariance $\Sigma_x$ has a small-eigenvalue tail.
The batch size scales as $B=\beta n$ and the vocabulary as $V=\alpha n^\gamma$.
Our main target is the \emph{width dependence} (as $n\to\infty$ with $\beta$ fixed) of the batch effect on $\eta^\star$.

The key technical issue is that the softmax loss is not quadratic. We give a self-contained derivation showing:
\begin{itemize}
\item under mild, checkable conditions the softmax one-step line search is governed (for \emph{scaling}) by a local quadratic surrogate,
\item in the coupled scaling $B=\beta n$ the SGD optimal step depends only on \emph{positive moments} of $\Sigma_x$,
\item whereas the Muon optimal step exhibits dependence on \emph{inverse-tail} quantities (e.g.\ $\tfrac1n\mathrm{Tr}(\Sigma_x^{-2})$) when the geometry is not in a special commuting case.
\end{itemize}

\section{Softmax model (with tied embeddings) and stochastic batches}

\subsection{Dimensions and scalings}
Let hidden width be $n$ and vocabulary size be $V=\alpha n^\gamma$ with constants $\alpha>0$ and $\gamma>0$.
Let batch size be $B=\beta n$ with $\beta\ge 1$.

\subsection{Representation distribution}
We model the pooled representation $x\in\mathbb{R}^n$ by a Gaussian ansatz
\[
x \sim \mathcal{N}(0,\Sigma_x),
\qquad \Sigma_x\succeq 0,
\]
where $\Sigma_x=\Sigma_x^{(n)}$ may depend on $n$ and has a lower spectral tail (specified later).
A batch $\{x_b\}_{b=1}^B$ yields the sample covariance
\[
S := \frac{1}{B}\sum_{b=1}^B x_b x_b^\top.
\]

\subsection{Parameters and logits}
Let $E\in\mathbb{R}^{n\times V}$ be a tied embedding matrix with iid Gaussian columns:
\[
E_{\cdot k}\sim \mathcal{N}(0,I_n)\ \text{iid across }k,
\]
independent of the batch.
We train a weight matrix $W\in\mathbb{R}^{n\times n}$ (we can imagine an operating point $\bar W$).
Given $x$, define logits
\[
z(W;x) := \frac{1}{n}E^\top W x \in \mathbb{R}^V,
\qquad
p(W;x) := \mathrm{softmax}(z(W;x)).
\]

\subsection{Labels and loss}
For a class label $t\in\{1,\dots,V\}$, we use \emph{label smoothing} with parameter $\varepsilon\in[0,1)$:
\[
y(t) := (1-\varepsilon)e_t + \varepsilon u,
\qquad
u:=\frac{1}{V}\mathbf{1}.
\]
The per-example loss is cross-entropy
\[
\ell(z,t):= -\sum_{k=1}^V y_k(t)\log p_k(z).
\]
Let $t$ be drawn from some distribution (e.g.\ Zipf); the batch sensitivity results below do not rely on the exact law of $t$
beyond mild non-saturation conditions at the operating point.

\section{One-step line search and the softmax Taylor expansion}

Fix an \emph{operating point} $\bar W$ and a \emph{batch-dependent update direction} $U_B$ (to be defined for SGD/Muon).
Consider the one-dimensional function
\[
f_B(\eta) := \mathbb{E}_{(x,t)}\big[\ell(z(\bar W-\eta U_B;x),t)\big],
\]
and define the one-step optimal learning rate
\[
\eta_B^\star := \arg\min_{\eta\ge 0} f_B(\eta).
\]

\subsection{Softmax derivatives in logit space}
Let $p=\mathrm{softmax}(z)$. Then
\[
g(z,t):=\nabla_z \ell(z,t)=p-y(t),
\qquad
H(z):=\nabla_z^2 \ell(z,t)=\mathrm{diag}(p)-pp^\top,
\]
where $H(z)$ is independent of $t$.

A basic global bound is $\|H(z)\|_{\mathrm{op}}\le \tfrac14$ for all $z$.

\subsection{Logit direction induced by a parameter direction}
Define the logit direction induced by $U_B$ as
\[
\Delta z(x) := \frac{1}{n}E^\top U_B x,
\qquad
z_\eta(x) := z(\bar W;x) - \eta\,\Delta z(x).
\]
Then $f_B'(\eta)= -\mathbb{E}\big[g(z_\eta,t)^\top \Delta z\big]$ and
\[
f_B''(\eta)= \mathbb{E}\big[\Delta z^\top H(z_\eta)\Delta z\big].
\]

\subsection{Second-order expansion and the cubic remainder}
Taylor expand $f_B$ at $\eta=0$:
\[
f_B(\eta)= f_B(0) - a_B \eta + \frac12 b_B \eta^2 + R_B(\eta),
\]
where
\[
a_B := -f_B'(0)=\mathbb{E}\big[g(z_0,t)^\top \Delta z\big],
\qquad
b_B := f_B''(0)=\mathbb{E}\big[\Delta z^\top H(z_0)\Delta z\big],
\quad z_0(x):=z(\bar W;x).
\]

A standard Taylor theorem bound expresses $R_B'(\eta)$ in terms of the third derivative of $\ell$ along the line.
For softmax cross-entropy, third derivatives w.r.t.\ logits are uniformly bounded in magnitude by an absolute constant,
hence there exists $C>0$ such that for all $\eta$ in any interval where the path stays non-saturated,
\begin{equation}
\label{eq:Rprime}
|R_B'(\eta)| \le C\,|\eta|^2\,\mathbb{E}\|\Delta z(x)\|_3^3,
\qquad
\|\Delta z\|_3^3:=\sum_{k=1}^V|\Delta z_k|^3.
\end{equation}

\section{Key lemma: $\mathbb{E}\|\Delta z\|_3^3$ scales like $\alpha n^{\gamma-3/2}$}

This is the quantitative ingredient that makes the softmax non-quadratic remainder negligible for $\gamma<3/2$,
including the main case $V\propto n$ ($\gamma=1$).

\subsection{A mild step-direction boundedness condition}
We assume the update direction is \emph{operator-norm bounded}:
\begin{equation}
\label{eq:UbBound}
\|U_B\|_{\mathrm{op}} = O(1)
\quad\text{uniformly in }n.
\end{equation}
This is the natural ``$\Theta(1)$ function-space step'' regime: combined with the $1/n$ factor in logits and $\|x\|=\Theta(\sqrt n)$,
it keeps each logit step coordinate at scale $n^{-1/2}$.

We also assume $\mathrm{Tr}(\Sigma_x)=\Theta(n)$, which holds in typical pooling/topic models.

\subsection{Coordinate scaling}
Condition on $(U_B,x)$. Since $E_{\cdot k}\sim\mathcal{N}(0,I_n)$ is independent of $(U_B,x)$,
\[
\Delta z_k \mid (U_B,x)\sim \mathcal{N}\!\left(0,\frac{1}{n^2}\|U_B x\|^2\right).
\]
Under \eqref{eq:UbBound} and $\mathbb{E}\|x\|^2=\mathrm{Tr}(\Sigma_x)=\Theta(n)$,
we have $\|U_B x\|^2\le \|U_B\|_{\mathrm{op}}^2\|x\|^2=O(n)$ and thus
\[
\mathbb{E}|\Delta z_k|^3 = \Theta(n^{-3/2}).
\]
Summing over $V=\alpha n^\gamma$ coordinates yields
\begin{equation}
\label{eq:l3scale}
\mathbb{E}\|\Delta z\|_3^3
=
\sum_{k=1}^V \mathbb{E}|\Delta z_k|^3
=
\Theta\!\left(\alpha\,n^{\gamma-3/2}\right).
\end{equation}

\subsection{Consequence for the softmax remainder}
Combining \eqref{eq:Rprime} and \eqref{eq:l3scale}, for $\gamma<3/2$ and $|\eta|=O(1)$ we obtain
\[
|R_B'(\eta)| = O\!\left(|\eta|^2\,\alpha\,n^{\gamma-3/2}\right)=o_n(1).
\]
In particular for $\gamma=1$, $|R_B'(\eta)|=O(\alpha n^{-1/2})$.

\section{From softmax to a quadratic surrogate with batch covariance $S$}

The second-order term $b_B$ admits an explicit quadratic form in $U_B$.

\subsection{Softmax-induced left factor $M$}
Let $\bar H:=\mathbb{E}[H(z_0(x))]$ (population-averaged logit Hessian at $\bar W$).
Define
\[
M := \frac{1}{n^2}E\,\bar H\,E^\top \in\mathbb{R}^{n\times n}.
\]
Then
\[
b_B = \mathbb{E}\big[\Delta z^\top H(z_0)\Delta z\big]
\approx \mathbb{E}\big[\Delta z^\top \bar H \Delta z\big]
= \mathbb{E}\big[x^\top U_B^\top M U_B x\big]
= \mathrm{Tr}\!\big(U_B^\top M U_B\,\Sigma_x\big),
\]
and on a batch, $\Sigma_x$ is replaced by $S$.

\subsection{The symmetry-breaking condition}
At strict symmetric initialization, $p\approx u$ implies $H$ is permutation-symmetric and $M\approx cI$, which kills the
non-commuting geometry needed for Muon sensitivity to small eigenvalues.
We therefore assume:
\begin{equation}
\label{eq:MnotI}
M \not\propto I
\quad\text{(equivalently $\bar H$ is not permutation-symmetric).}
\end{equation}
A simple way to motivate \eqref{eq:MnotI} without ``running training'' is to take $\bar W=W_0+\epsilon W_1$ with small $\epsilon$
and data-correlated $W_1$, so that predictions are slightly non-uniform and $\bar H$ is not symmetric across classes.

\subsection{Quadratic proxy for one-step scaling}
Under non-saturation and $\gamma<3/2$, the remainder derivative is $o_n(1)$ for $|\eta|=O(1)$, hence
the location of the minimizer (for scaling with $n$) is governed by the quadratic approximation
\[
f_B(\eta) \approx f_B(0) - a_B \eta + \frac12 b_B \eta^2,
\qquad \eta^\star_{B,\mathrm{quad}}=\frac{a_B}{b_B}.
\]
Non-quadratic softmax effects shift $\eta^\star$ by at most $O(\alpha n^{\gamma-3/2})$ under mild convexity along the line,
and therefore cannot create $\log n$ or power-law divergences when $\gamma=1$.

From this point onward, batch dependence enters \emph{only} through the random covariance $S$,
matching the quadratic/Wishart setup.

\section{Quadratic/Wishart setup for optimizer comparison (SGD vs Muon)}

To isolate the optimizer effect, consider the canonical quadratic objective in an ``error'' variable $E\in\mathbb{R}^{n\times n}$:
\[
L(E)=\frac12\|E\|_F^2.
\]
Given a batch covariance $S$, the minibatch gradient is
\[
g(E;S)=ES.
\]
This is the same algebraic structure as the softmax local quadratic term once we absorb left factors (e.g.\ $M^{1/2}$)
into the definition of the error coordinates. The spectrum and tail behavior enter through $S$ (hence $\Sigma_x$).

Let $A:=E^\top E\succeq 0$.

\subsection{SGD: exact one-step optimal learning rate}
SGD update: $E^+=E-\eta ES$.
A direct expansion gives
\[
L(E^+)=\frac12\mathrm{Tr}(A) - \eta\,\mathrm{Tr}(AS) + \frac{\eta^2}{2}\mathrm{Tr}(SAS).
\]
Thus the (conditional) one-step optimal learning rate is
\[
\eta_{\mathrm{opt}}^{\mathrm{SGD}}(E,S)=\frac{\mathrm{Tr}(AS)}{\mathrm{Tr}(SAS)}.
\]
Taking expectation over $S$ and using the Wishart moment identity for Gaussian $x$ yields
\[
\mathbb{E}[\mathrm{Tr}(SAS)]
=
\mathrm{Tr}(\Sigma_x A \Sigma_x)
+
\frac{1}{B}
\Big(
\mathrm{Tr}(\Sigma_x)\mathrm{Tr}(A\Sigma_x)
+\mathrm{Tr}(A\Sigma_x^2)
\Big).
\]
With $B=\beta n$, the correction scales like $1/\beta$ times \emph{positive} spectral moments of $\Sigma_x$.
In particular small eigenvalues of $\Sigma_x$ do not cause divergences in these terms.

\subsection{Muon: exact one-step optimal learning rate and inverse-tail sensitivity}
Muon update:
\[
Q=\mathrm{polar}(ES),
\qquad
E^+=E-\eta Q.
\]
Because $Q$ has orthonormal columns on the range of $ES$, one obtains the exact conditional optimizer
\[
\eta_{\mathrm{opt}}^{\mathrm{Muon}}(E,S)
=\frac{1}{n}\mathrm{Tr}(E^\top Q)
=\frac{1}{n}\mathrm{Tr}\!\left(S^{-1}(SAS)^{1/2}\right).
\]
Unlike SGD, this depends on $S$ through a \emph{matrix inverse} and a \emph{square root}, which makes it sensitive to
eigenvector perturbations of $S$ and, in generic non-commuting cases, to inverse-tail functionals.

A standard perturbative expansion of the polar/square-root functional around $\Sigma_x$ yields, for large $B$,
\[
\eta_{\mathrm{opt}}^{\mathrm{Muon}}(\beta,n)
=
\eta_\infty^{\mathrm{Muon}} - \frac{1}{\beta}\,C(A,\Sigma_x) + o(1/\beta),
\]
where $\eta_\infty^{\mathrm{Muon}}=\tfrac{1}{n}\mathrm{Tr}\big(\Sigma_x^{-1}(\Sigma_x A \Sigma_x)^{1/2}\big)$ and
$C(A,\Sigma_x)$ is an eigenvector-perturbation coefficient.
A representative form (in the eigenbasis of $\Sigma_x$) is
\[
C(A,\Sigma_x)\ \approx\ \frac{1}{n}\sum_{i\ne j}\frac{\lambda_i\lambda_j}{(\lambda_i-\lambda_j)^2}\big(\sqrt{a_i}-\sqrt{a_j}\big)^2,
\]
where $\{\lambda_i\}$ are eigenvalues of $\Sigma_x$ and $\{a_i\}$ are eigenvalues of the teacher-induced matrix in that basis
(in commuting cases, the numerator vanishes and $C=0$).

\paragraph{Generic non-commuting scaling.}
In generic non-commuting situations (no shared eigenbasis and no special cancellations), the dominant scaling is controlled by inverse-tail quantities,
schematically
\[
C(A,\Sigma_x)\ \asymp\ \kappa(A,\text{noncomm})\cdot \frac{1}{n}\mathrm{Tr}(\Sigma_x^{-2}),
\]
where $\kappa(A,\text{noncomm})$ measures the strength of non-commutativity (and $\kappa=0$ in special commuting/spherical cases).
This is the mechanism by which Muon becomes sensitive to small eigenvalues of $\Sigma_x$.

\section{Width scaling with $B=\beta n$ under a power-law lower tail}

Assume the eigenvalues of $\Sigma_x=\Sigma_x^{(n)}$ satisfy a power-law to zero:
\[
\lambda_i(\Sigma_x)\asymp (i/n)^\alpha,\qquad i=1,\dots,n,
\]
so $\lambda_{\min}\asymp n^{-\alpha}$.
Then inverse moments scale as
\[
\frac{1}{n}\mathrm{Tr}(\Sigma_x^{-2})
\sim
\begin{cases}
O(1), & \alpha<\tfrac12,\\[4pt]
\Theta(\log n), & \alpha=\tfrac12,\\[4pt]
\Theta(n^{2\alpha-1}), & \alpha>\tfrac12.
\end{cases}
\]
Therefore, in the coupled regime $B=\beta n$, the Muon one-step optimizer exhibits a width-dependent batch effect
\[
\eta_{\mathrm{opt}}^{\mathrm{Muon}}(\beta,n)
=
\eta_\infty^{\mathrm{Muon}}
-
\frac{\kappa}{\beta}\times
\begin{cases}
O(1),\\
\Theta(\log n),\\
\Theta(n^{2\alpha-1}),
\end{cases}
\]
while the SGD optimizer depends on $\Sigma_x$ only through positive moments and has no analogous inverse-tail divergence.

\section{Putting it back into the softmax model}

Under the softmax-to-quadratic reduction:
\begin{itemize}
\item $\Sigma_x$ is the covariance of pooled representations, which can naturally have a power-law lower tail (e.g.\ from topic mixtures and pooling).
\item The softmax-induced left factor is $M=\tfrac1{n^2}E\bar H E^\top$; to avoid the symmetric special case, assume $M\not\propto I$.
\item For $V=\alpha n^\gamma$ with $\gamma<3/2$ (in particular $\gamma=1$), the softmax non-quadratic remainder along one step obeys
$|R_B'(\eta)|=O(\alpha n^{\gamma-3/2})$, hence it cannot produce the $\log n$ or power-law width dependence that arises from inverse-tail
effects of $\Sigma_x$ in the Muon polar map.
\end{itemize}

\paragraph{Interpretation.}
In the width--batch coupled scaling $B=\beta n$, finite-batch randomness persists at the level of eigen-structure of $S$,
and Muon couples to this through polar/square-root operations that are amplified by small eigenvalues unless the geometry is in a commuting/spherical special case.
SGD couples only through polynomial (positive-moment) functionals of $S$ and remains stable to lower-tail behavior.

\section{Summary of assumptions (minimal for the scaling conclusions)}
\begin{enumerate}
\item $V=\alpha n^\gamma$ with $\gamma<3/2$ (main case $\gamma=1$), and $B=\beta n$ with $\beta\ge 1$.
\item $x\sim\mathcal N(0,\Sigma_x^{(n)})$ with $\mathrm{Tr}(\Sigma_x)=\Theta(n)$ and a specified lower spectral tail.
\item Update directions satisfy $\|U_B\|_{\mathrm{op}}=O(1)$ so per-coordinate logit steps are $O(n^{-1/2})$.
\item Non-saturation along the relevant $\eta$ range (helped by label smoothing $\varepsilon>0$), so $f_B''$ stays $\Theta(1)$ in the normalized regime.
\item Symmetry breaking at the operating point: $M=\tfrac1{n^2}E\bar H E^\top \not\propto I$.
\item Generic non-commuting geometry for Muon to exhibit inverse-tail scaling; commuting/spherical cases yield $\kappa=0$ and suppress the effect.
\end{enumerate}


\section{Setup}

Let $x \sim \mathcal{N}(0,\Sigma)$ with $\Sigma \succeq 0$.
Let the batch size be $B = c n$ with $c \ge 1$.
Define the sample covariance
\[
S = \frac{1}{B} \sum_{b=1}^B x_b x_b^\top.
\]

We study the quadratic objective
\[
L(E) = \frac{1}{2}\|E\|_F^2,
\qquad E = W - W_\star,
\]
so the minibatch gradient is
\[
g = ES.
\]

We analyze the \emph{first optimization step} and compute the one-step optimal
learning rate for both SGD and Muon.

\section{SGD: Exact First-Step Optimal Learning Rate}

SGD update:
\[
E^+ = E - \eta ES.
\]

Let $A = E^\top E$. Then
\[
L(E^+) = \frac{1}{2} \mathrm{Tr}(A)
- \eta \mathrm{Tr}(AS)
+ \frac{\eta^2}{2} \mathrm{Tr}(SAS).
\]

Taking expectation over the batch,
\[
\eta_{\mathrm{opt}}^{\mathrm{SGD}}(E)
= \frac{\mathrm{Tr}(A\Sigma)}
{\mathbb{E}[\mathrm{Tr}(SAS)]}.
\]

For Gaussian data,
\[
\mathbb{E}[\mathrm{Tr}(SAS)]
=
\mathrm{Tr}(\Sigma A \Sigma)
+
\frac{1}{B}
\left(
\mathrm{Tr}(\Sigma)\mathrm{Tr}(A\Sigma)
+
\mathrm{Tr}(A\Sigma^2)
\right).
\]

Thus, with $B=cn$,
\[
\eta_\infty^{\mathrm{SGD}}(c)
=
\frac{\mathrm{Tr}(A\Sigma)}
{\mathrm{Tr}(\Sigma A \Sigma)}
\left(
1
-
\frac{1}{c}
\frac{
\mathrm{Tr}(\Sigma)\mathrm{Tr}(A\Sigma)
+
\mathrm{Tr}(A\Sigma^2)
}
{\mathrm{Tr}(\Sigma A \Sigma)}
+ O(c^{-2})
\right).
\]

\subsection*{Key Observation}

The $1/c$ correction depends only on \emph{positive moments}
of $\Sigma$ (e.g.\ $\mathrm{Tr}(\Sigma)$ and $\mathrm{Tr}(\Sigma^2)$).
A lower spectral tail near $0$ does not produce divergences.

\section{Muon: Exact First-Step Optimal Learning Rate}

Muon update:
\[
Q = \mathrm{polar}(ES),
\qquad
E^+ = E - \eta Q.
\]

The one-step optimal learning rate is
\[
\eta_{\mathrm{opt}}^{\mathrm{Muon}}(E,S)
=
\frac{1}{n}
\mathrm{Tr}
\left(
E^\top \mathrm{polar}(ES)
\right).
\]

Equivalently,
\[
\eta_{\mathrm{opt}}^{\mathrm{Muon}}(E,S)
=
\frac{1}{n}
\mathrm{Tr}
\left(
S^{-1}(SAS)^{1/2}
\right).
\]

For large $c$,
\[
\eta_\infty^{\mathrm{Muon}}(c)
=
\eta_\infty^{\mathrm{Muon}}(\infty)
-
\frac{C(A,\Sigma)}{c}
+ o(c^{-1}),
\]
where
\[
\eta_\infty^{\mathrm{Muon}}(\infty)
=
\frac{1}{n}
\mathrm{Tr}
\left(
\Sigma^{-1}
(\Sigma A \Sigma)^{1/2}
\right).
\]

\section{Power-Law Spectrum Extending to Zero}

Assume the eigenvalues of $\Sigma$ satisfy
\[
\lambda_i(\Sigma) \asymp (i/n)^\alpha,
\qquad i=1,\dots,n,
\]
so $\lambda_{\min} \asymp n^{-\alpha}$.

Inverse moments scale as
\[
\frac{1}{n}\mathrm{Tr}(\Sigma^{-2})
\sim
\begin{cases}
O(1), & \alpha < \tfrac12,\\
\Theta(\log n), & \alpha = \tfrac12,\\
\Theta(n^{2\alpha-1}), & \alpha > \tfrac12.
\end{cases}
\]

\section{Teacher Structure Cases}

Let $A = E_0^\top E_0$ be determined by the teacher.

\subsection*{Case A: Spherical Teacher}

If $A = \alpha I$, then
\[
(SAS)^{1/2} = \sqrt{\alpha}\, S,
\qquad
S^{-1}(SAS)^{1/2} = \sqrt{\alpha}\, I.
\]

Therefore,
\[
\eta_{\mathrm{opt}}^{\mathrm{Muon}} = \sqrt{\alpha},
\]
independent of $c$ and $\Sigma$.
Thus
\[
C(A,\Sigma) = 0.
\]

\subsection*{Case B: Teacher Aligned with $\Sigma$}

Let $A = \mathrm{diag}(a_i)$ in the eigenbasis of $\Sigma$.
Then for large $c$,
\[
\eta_\infty^{\mathrm{Muon}}(\infty)
=
\frac{1}{n}\sum_{i=1}^n \sqrt{a_i}.
\]

The correction coefficient behaves like
\[
C(A,\Sigma)
\approx
\frac{1}{n}
\sum_{i \ne j}
\frac{\lambda_i \lambda_j}{(\lambda_i - \lambda_j)^2}
\big(\sqrt{a_i} - \sqrt{a_j}\big)^2.
\]

\subsection*{Smooth Teacher Spectrum}

If $\sqrt{a_i} = f(\lambda_i)$ with $f$ Lipschitz near $0$,
then
\[
C(A,\Sigma)
\sim
\frac{1}{n}
\sum_i
\lambda_i^2 f'(\lambda_i)^2,
\]
which remains finite under mild regularity.

\subsection*{Tail-Amplifying Teacher}

If $\sqrt{a_i} = \lambda_i^{-\beta}$,
then
\[
C(A,\Sigma)
\propto
\int_0
\lambda^{1/\alpha - 1 - 2\beta}
\, d\lambda,
\]
which diverges when
\[
\beta \ge \frac{1}{2\alpha}.
\]

\subsection*{Generic Non-Aligned Teacher}

If $A$ does not commute with $\Sigma$,
\[
C(A,\Sigma)
\sim
(\text{non-commutativity of }A\text{ and }\Sigma)
\times
\frac{1}{n}\mathrm{Tr}(\Sigma^{-2}).
\]

Thus for $\alpha \ge \tfrac12$,
\[
C(A,\Sigma)
=
\begin{cases}
O(1), & \alpha < \tfrac12,\\
O(\log n), & \alpha = \tfrac12,\\
O(n^{2\alpha-1}), & \alpha > \tfrac12.
\end{cases}
\]

\section{Conclusion}

\begin{itemize}
\item SGD batch corrections depend only on positive moments of $\Sigma$ and remain stable under heavy lower spectral tails.
\item Muon batch corrections depend on eigenvector perturbations and inverse-tail quantities such as $\mathrm{Tr}(\Sigma^{-2})$.
\item For power-law spectra extending to zero, Muon becomes increasingly sensitive to tail behavior unless the teacher induces commuting/spherical structure.
\end{itemize}

\end{document}
